\usepackage[utf8]{inputenc}
\usepackage[T1]{fontenc}
\usepackage[english]{babel}
% Provadis "Anforderungen" §3.2.4 prescribes Times New Roman 12, Arial 11 or
% Verdana 10 for the running text. mathptmx supplies a Times clone for text and
% maths, helvet the matching sans. The typewriter face deliberately stays Latin
% Modern rather than Courier: Courier is far wider and would overflow the narrow
% feature-reference table columns that carry long snake_case identifiers.
\usepackage{mathptmx}
\usepackage[scaled=0.90]{helvet}
\renewcommand{\ttdefault}{lmtt}
\usepackage{microtype}
\usepackage{csquotes}
\usepackage{geometry}
\usepackage{setspace}
\usepackage{graphicx}
\usepackage{booktabs}
\usepackage{amsmath}
\usepackage{amssymb}
\usepackage{tabularx}
\usepackage{longtable}
\usepackage{enumitem}
\usepackage{xcolor}
\usepackage{listings}
\usepackage{caption}
\usepackage{subcaption}
\usepackage{tikz}
\usetikzlibrary{arrows.meta,decorations.pathreplacing,calc,backgrounds}
\usepackage[nohyperlinks]{acronym}
% Used to bind the signed Telekom Steckbrief into the front matter (see main.tex).
\usepackage{pdfpages}

% Reference from the running text into the KI-Verzeichnis, as required by §5.1 of the
% Provadis "Regelungen zum Wissenschaftlichen Arbeiten mit Hilfe von Künstlicher
% Intelligenz". \kiref{Claude Opus 5/3} marks verbatim adoption, \kirefcf{...} marks a
% passage that was only developed from the AI output ("vgl.").
\newcommand{\kiref}[1]{\footnote{#1 --- see the Declaration on the Use of Artificial
  Intelligence, p.~\pageref{chap:ai-usage}.}}
\newcommand{\kirefcf}[1]{\footnote{cf.\ #1 --- see the Declaration on the Use of
  Artificial Intelligence, p.~\pageref{chap:ai-usage}.}}

% Loud placeholder for the verbatim opening prompts that still have to be pasted into
% the KI-Verzeichnis. Impossible to miss in the compiled PDF.
\newcommand{\promptTODO}[1]{\textbf{[\,PROMPT #1 --- TO BE INSERTED\,]}}

% The back matter reaches seven-character Roman page numbers (e.g. XXXVIII), which
% overflow the default 1.55em page-number field of the contents lists. Affects the
% lists only; margins, typeface and page numbering itself are unchanged.
\makeatletter
\renewcommand{\@pnumwidth}{4em}
\renewcommand{\@tocrmarg}{5em}
\makeatother
\usepackage[
  backend=biber,
  style=ieee,
  sorting=none
]{biblatex}
% Long bibliography URLs are otherwise unbreakable and run into the right margin.
% xurl permits a break after any character in a URL; the biburl penalties make biblatex
% actually take those breaks rather than only allowing them at slashes.
\usepackage{xurl}
\setcounter{biburllcpenalty}{7000}
\setcounter{biburlucpenalty}{8000}
\setcounter{biburlnumpenalty}{7000}

\usepackage[hidelinks]{hyperref}
\usepackage[nameinlink,noabbrev]{cleveref}

% Provadis "Anforderungen" §3.2.4: links und oben je 3cm, rechts und unten je 2cm.
\geometry{
  left=30mm,
  right=20mm,
  top=30mm,
  bottom=20mm
}

\onehalfspacing
\setlength{\parindent}{0pt}
\setlength{\parskip}{0.7em}

% Heading sizes and the space above a chapter title.
%
% §3.2.4 of the Anforderungen recommends Times New Roman 14, Arial 13 or Verdana 12
% for chapter headings. Headings are set in the helvet sans (the Arial equivalent),
% so the applicable figure is 13pt. Because helvet is loaded with scaled=0.90, a
% nominal size renders at 0.90 of its value, hence 14.4pt nominal for 13.0pt
% effective. The lower levels step down from there so the hierarchy is preserved;
% \sffamily is repeated because \setkomafont replaces the whole declaration and
% would otherwise fall back to the serif body face.
\setkomafont{chapter}{\sffamily\bfseries\fontsize{14.4}{17.3}\selectfont}
\setkomafont{section}{\sffamily\bfseries\fontsize{13.2}{15.8}\selectfont}
\setkomafont{subsection}{\sffamily\bfseries\fontsize{12}{14.4}\selectfont}
%
% scrreprt defaults the space above a chapter title to 3.3\baselineskip+\parskip.
% The Anforderungen say nothing about it, and with \onehalfspacing the default
% leaves roughly 2.8cm of white space below an already 3cm top margin.
\RedeclareSectionCommand[beforeskip=-1.5\baselineskip,afterskip=.8\baselineskip]{chapter}

% Breakable underscore for long snake_case feature identifiers in narrow table cells.
\newcommand{\su}{\_\allowbreak}
% Typeset a feature key in monospace (underscores written as \su so the key can wrap).
\newcommand{\feat}[1]{\texttt{#1}}

% Long snake_case identifiers in running prose cannot be broken. Permit a line break
% after every underscore, as \su already does inside the feature tables. \allowbreak is
% not expandable into a PDF string, so it is neutralized for hyperref bookmarks.
\let\origunderscore\_
\renewcommand{\_}{\origunderscore\allowbreak}
\pdfstringdefDisableCommands{\let\allowbreak\empty}

% \texttt cannot break at '.', '-' or '/', so file paths and hyphenated identifiers run
% into the margin. These mirror \su: the character, followed by a permitted break.
\newcommand{\sdot}{.\allowbreak}
\newcommand{\sdash}{-\allowbreak}
\newcommand{\sslash}{/\allowbreak}

% With those breakpoints available, let TeX stretch interword space a little rather than
% choosing a slightly overfull line over a loose one.
\setlength{\emergencystretch}{1.5em}

\addbibresource{references.bib}

\graphicspath{{images/}}

\lstdefinestyle{thesiscode}{
  basicstyle=\ttfamily\small,
  breaklines=true,
  frame=single,
  numbers=left,
  numberstyle=\tiny,
  tabsize=2,
  showstringspaces=false
}
\lstdefinelanguage{JavaScript}{
  keywords={break,const,continue,else,for,function,if,import,let,return,var,while},
  sensitive=true,
  comment=[l]{//},
  morecomment=[s]{/*}{*/},
  morestring=[b]",
  morestring=[b]'
}
\lstset{style=thesiscode}
